\section{Funciones del Producto}

Esta sección especifica los requisitos funcionales del framework Kit Pagos Colombia. La sección 5.1 presenta los requisitos funcionales en formato tabular con identificador, descripción, componente y prioridad. La sección 5.2 en adelante complementa esta especificación mediante historias de usuario que expresan cada requisito desde la perspectiva del usuario, acompañadas de criterios de aceptación verificables.

\subsection{Requisitos Funcionales}

La Tabla 5 presenta el listado consolidado de requisitos funcionales del framework, organizados por componente. Las prioridades se clasifican en Alta (esencial para el MVP), Media (importante pero no bloqueante) y Baja (deseable en versiones futuras).

\begin{center}
\renewcommand{\arraystretch}{1.2}
\begin{longtable}{|c|p{8.5cm}|c|c|}
  \hline
  \textbf{ID} & \textbf{Descripción del requisito} & \textbf{Componente} & \textbf{Prioridad} \\
  \hline
  \endfirsthead
  \hline
  \textbf{ID} & \textbf{Descripción del requisito} & \textbf{Componente} & \textbf{Prioridad} \\
  \hline
  \endhead
  \hline
  \endfoot
  \hline
  \caption{Requisitos funcionales del framework Kit Pagos Colombia.}
  \label{tab:requisitos-funcionales}\\
  \endlastfoot
  RF-01 & El sistema debe permitir crear una intención de pago especificando pasarela, monto, moneda, referencia de orden y datos del pagador mediante una interfaz unificada. & SDK & Alta \\
  \hline
  RF-02 & El sistema debe retornar una respuesta normalizada ante la creación de un pago, incluyendo el identificador de transacción, el estado inicial y la URL o token de pago según la pasarela. & SDK & Alta \\
  \hline
  RF-03 & El sistema debe permitir consultar el estado de una transacción por identificador y retornarlo mapeado a los estados normalizados: PENDIENTE, APROBADO, RECHAZADO, EXPIRADO o ERROR. & SDK & Alta \\
  \hline
  RF-04 & El sistema debe validar la firma criptográfica de un webhook entrante utilizando el algoritmo correspondiente a cada pasarela y retornar un evento normalizado si la firma es válida. & SDK & Alta \\
  \hline
  RF-05 & El sistema debe lanzar excepciones tipadas con código de error normalizado y mensaje descriptivo ante cualquier condición de fallo esperada (timeout, credenciales inválidas, recurso no encontrado). & SDK & Alta \\
  \hline
  RF-06 & El sistema debe permitir cambiar la pasarela activa modificando únicamente el parámetro de configuración en la inicialización del framework, sin modificar el código de negocio del llamador. & SDK & Alta \\
  \hline
  RF-07 & El sistema debe reintentar automáticamente una operación ante errores transitorios (timeout, respuestas 5xx) siguiendo una estrategia de backoff exponencial configurable, hasta un número máximo de intentos definido por el usuario. & SDK & Media \\
  \hline
  RF-08 & El sistema debe garantizar que las credenciales de cada pasarela no sean expuestas en logs, mensajes de error ni en las respuestas del SDK. & SDK & Alta \\
  \hline
  RF-09 & El sistema debe exponer endpoints REST que repliquen el comportamiento de creación de pago de cada pasarela soportada, accesibles sin credenciales reales. & API Simulación & Alta \\
  \hline
  RF-10 & El sistema debe permitir configurar el escenario de respuesta de cada endpoint mediante un parámetro de la solicitud, soportando al menos: APROBADO, RECHAZADO, FONDOS\_INSUFICIENTES, TIMEOUT y ERROR\_RED. & API Simulación & Alta \\
  \hline
  RF-11 & El sistema debe replicar la estructura exacta de las respuestas de error de cada pasarela en los escenarios de fallo configurados. & API Simulación & Alta \\
  \hline
  RF-12 & El sistema debe permitir solicitar el envío de un webhook simulado hacia una URL destino, con el formato y firma correspondientes a la pasarela indicada. & API Simulación & Media \\
  \hline
  RF-13 & El sistema debe exponer la especificación OpenAPI 3.0 de la API de simulación en un endpoint público (/docs), accesible a través de un navegador web. & API Simulación & Media \\
  \hline
  RF-14 & El sistema debe proveer documentación de datos de prueba organizada por pasarela, incluyendo tarjetas de prueba, credenciales de sandbox y referencias de pago válidas para cada proveedor. & Documentación & Alta \\
  \hline
  RF-15 & El sistema debe indicar el comportamiento esperado de cada dato de prueba (tarjeta aprobada, tarjeta rechazada, fondos insuficientes) junto con la fuente oficial del proveedor. & Documentación & Media \\
  \hline
\end{longtable}
\end{center}

\subsection{Historias de Usuario}

Las siguientes historias de usuario complementan los requisitos funcionales de la Tabla 5, expresando cada necesidad desde la perspectiva del usuario y acompañándola de criterios de aceptación verificables. Cada historia referencia el requisito funcional que desarrolla.

\subsection{SDK — Interfaz unificada de integración}

\subsubsection{HU-01: Crear una intención de pago con una pasarela específica}
\begin{table}[htbp]
  \centering
  \renewcommand{\arraystretch}{1.2}
  \begin{tabularx}{\textwidth}{|l|X|}
    \hline
    \textbf{Como} & desarrollador de software \\
    \hline
    \textbf{Quiero} & crear una intención de pago especificando únicamente la pasarela deseada, el monto y los datos del pagador mediante una sola llamada al SDK \\
    \hline
    \textbf{Para} & no tener que conocer los detalles de autenticación y formato de solicitud propios de cada proveedor \\
    \hline
  \end{tabularx}
  \caption{HU-01}
  \label{tab:hu-01}
\end{table}

\textbf{Criterios de aceptación:}
\begin{itemize}
  \item El SDK acepta una solicitud de creación de pago con los parámetros: pasarela, monto, moneda, referencia de orden y datos del pagador.
  \item El SDK retorna una respuesta normalizada con el identificador de la transacción, el estado inicial y la URL de redirección o token de pago según corresponda.
  \item Si la pasarela retorna un error, el SDK lanza una excepción tipada con código de error normalizado y mensaje descriptivo.
  \item El SDK soporta las pasarelas Wompi, PayU, Mercado Pago y Kushki sin que el código del llamador deba cambiar entre ellas.
\end{itemize}

\subsubsection{HU-02: Consultar el estado de una transacción}
\begin{table}[htbp]
  \centering
  \renewcommand{\arraystretch}{1.2}
  \begin{tabularx}{\textwidth}{|l|X|}
    \hline
    \textbf{Como} & desarrollador de software \\
    \hline
    \textbf{Quiero} & consultar el estado actual de una transacción utilizando su identificador, independientemente de la pasarela que la procesó \\
    \hline
    \textbf{Para} & obtener siempre un estado normalizado sin tener que interpretar los estados específicos de cada proveedor \\
    \hline
  \end{tabularx}
  \caption{HU-02}
  \label{tab:hu-02}
\end{table}

\textbf{Criterios de aceptación:}
\begin{itemize}
  \item El SDK retorna el estado de la transacción mapeado a uno de los estados normalizados: PENDIENTE, APROBADO, RECHAZADO, EXPIRADO o ERROR.
  \item La respuesta incluye el estado original de la pasarela como campo adicional para referencia.
  \item Si el identificador no existe en la pasarela indicada, el SDK lanza una excepción tipada con código NOT\_FOUND.
\end{itemize}

\subsubsection{HU-03: Validar y procesar un webhook entrante}
\begin{table}[htbp]
  \centering
  \renewcommand{\arraystretch}{1.2}
  \begin{tabularx}{\textwidth}{|l|X|}
    \hline
    \textbf{Como} & desarrollador de software \\
    \hline
    \textbf{Quiero} & validar la firma de un webhook recibido y extraer el evento normalizado que contiene, con una sola llamada al SDK \\
    \hline
    \textbf{Para} & no tener que implementar la lógica de validación de firma de cada pasarela por separado \\
    \hline
  \end{tabularx}
  \caption{HU-03}
  \label{tab:hu-03}
\end{table}

\textbf{Criterios de aceptación:}
\begin{itemize}
  \item El SDK valida la firma del webhook utilizando el algoritmo correspondiente a cada pasarela (HMAC-SHA256 para Wompi y Kushki, firma propia de PayU y Mercado Pago).
  \item Si la firma es inválida, el SDK retorna un error tipado con código INVALID\_SIGNATURE.
  \item Si la firma es válida, el SDK retorna un objeto de evento normalizado con el tipo de evento, el identificador de transacción y el nuevo estado.
\end{itemize}

\subsubsection{HU-04: Cambiar de pasarela sin modificar la lógica de negocio}
\begin{table}[htbp]
  \centering
  \renewcommand{\arraystretch}{1.2}
  \begin{tabularx}{\textwidth}{|l|X|}
    \hline
    \textbf{Como} & desarrollador de software senior \\
    \hline
    \textbf{Quiero} & cambiar la pasarela activa de mi aplicación modificando únicamente la configuración del framework, sin tocar el código de negocio \\
    \hline
    \textbf{Para} & reducir el riesgo y el esfuerzo de migración entre proveedores \\
    \hline
  \end{tabularx}
  \caption{HU-04}
  \label{tab:hu-04}
\end{table}

\textbf{Criterios de aceptación:}
\begin{itemize}
  \item El framework permite especificar la pasarela activa mediante un parámetro de configuración en la inicialización.
  \item Al cambiar la pasarela en la configuración, todas las operaciones de pago se redirigen al nuevo proveedor sin modificar el código del llamador.
  \item El cambio de pasarela no afecta la estructura ni el tipo de las respuestas retornadas por el SDK.
\end{itemize}

\subsubsection{HU-05: Manejo automático de reintentos ante fallos transitorios}
\begin{table}[htbp]
  \centering
  \renewcommand{\arraystretch}{1.2}
  \begin{tabularx}{\textwidth}{|l|X|}
    \hline
    \textbf{Como} & desarrollador de software \\
    \hline
    \textbf{Quiero} & que el SDK reintente automáticamente una operación cuando la pasarela devuelve un error transitorio (timeout, error 500) \\
    \hline
    \textbf{Para} & mejorar la resiliencia de la integración sin implementar lógica de reintento manualmente \\
    \hline
  \end{tabularx}
  \caption{HU-05}
  \label{tab:hu-05}
\end{table}

\textbf{Criterios de aceptación:}
\begin{itemize}
  \item El SDK reintenta la operación hasta un máximo configurable de veces ante errores de red o respuestas 5xx de la pasarela.
  \item Si se supera el número máximo de reintentos, el SDK lanza una excepción tipada con código MAX\_RETRIES\_EXCEEDED.
\end{itemize}

\subsection{API web de simulación}

\subsubsection{HU-06: Simular una transacción aprobada}
\begin{table}[htbp]
  \centering
  \renewcommand{\arraystretch}{1.2}
  \begin{tabularx}{\textwidth}{|l|X|}
    \hline
    \textbf{Como} & desarrollador o tester \\
    \hline
    \textbf{Quiero} & enviar una solicitud de pago a la API de simulación y recibir una respuesta que replica el comportamiento de una transacción aprobada por cada pasarela \\
    \hline
    \textbf{Para} & probar el flujo de éxito de mi integración sin depender del sandbox real de la pasarela \\
    \hline
  \end{tabularx}
  \caption{HU-06}
  \label{tab:hu-06}
\end{table}

\textbf{Criterios de aceptación:}
\begin{itemize}
  \item La API de simulación expone un endpoint por pasarela que acepta los mismos parámetros que la API real del proveedor.
  \item La respuesta del endpoint de éxito replica la estructura exacta de una respuesta aprobada de la pasarela correspondiente.
  \item La API de simulación es accesible públicamente sin requerir credenciales reales ni configuración local.
\end{itemize}

\subsubsection{HU-07: Simular escenarios de error controlados}
\begin{table}[htbp]
  \centering
  \renewcommand{\arraystretch}{1.2}
  \begin{tabularx}{\textwidth}{|l|X|}
    \hline
    \textbf{Como} & tester de software \\
    \hline
    \textbf{Quiero} & provocar escenarios de rechazo, timeout y error de red en la API de simulación mediante parámetros de la solicitud \\
    \hline
    \textbf{Para} & verificar que mi integración maneja correctamente los casos de fallo sin tener que reproducirlos en el sandbox real \\
    \hline
  \end{tabularx}
  \caption{HU-07}
  \label{tab:hu-07}
\end{table}

\textbf{Criterios de aceptación:}
\begin{itemize}
  \item La API de simulación soporta al menos los siguientes escenarios configurables por parámetro: APROBADO, RECHAZADO, FONDOS\_INSUFICIENTES, TIMEOUT, ERROR\_RED.
  \item La respuesta de cada escenario replica la estructura de error real de la pasarela correspondiente.
  \item El escenario TIMEOUT introduce un retardo configurable antes de responder, simulando latencia real.
\end{itemize}

\subsubsection{HU-08: Recibir una notificación webhook simulada}
\begin{table}[htbp]
  \centering
  \renewcommand{\arraystretch}{1.2}
  \begin{tabularx}{\textwidth}{|l|X|}
    \hline
    \textbf{Como} & desarrollador de software \\
    \hline
    \textbf{Quiero} & solicitar a la API de simulación que envíe un webhook a una URL de mi aplicación, replicando el formato y firma de la pasarela correspondiente \\
    \hline
    \textbf{Para} & probar el endpoint receptor de webhooks de mi aplicación sin esperar eventos reales del sandbox \\
    \hline
  \end{tabularx}
  \caption{HU-08}
  \label{tab:hu-08}
\end{table}

\textbf{Criterios de aceptación:}
\begin{itemize}
  \item La API de simulación acepta una URL destino y un tipo de evento, y envía un webhook firmado replicando el formato de la pasarela indicada.
  \item La firma del webhook simulado es válida y puede ser verificada con el mecanismo de validación del SDK.
  \item La API de simulación retorna el resultado del intento de envío (éxito o fallo de la entrega).
\end{itemize}

\subsection{Documentación de datos de prueba}

\subsubsection{HU-09: Acceder a datos de prueba por pasarela}
\begin{table}[htbp]
  \centering
  \renewcommand{\arraystretch}{1.2}
  \begin{tabularx}{\textwidth}{|l|X|}
    \hline
    \textbf{Como} & desarrollador o tester \\
    \hline
    \textbf{Quiero} & encontrar en un único lugar los datos de prueba disponibles para cada pasarela: tarjetas, credenciales de sandbox, referencias y parámetros especiales \\
    \hline
    \textbf{Para} & no tener que buscar en múltiples fuentes de documentación de cada proveedor durante el desarrollo \\
    \hline
  \end{tabularx}
  \caption{HU-09}
  \label{tab:hu-09}
\end{table}

\textbf{Criterios de aceptación:}
\begin{itemize}
  \item El repositorio del framework incluye una sección de documentación de datos de prueba organizada por pasarela.
  \item Para cada pasarela se documentan: tarjetas de prueba con sus comportamientos esperados, credenciales de acceso al sandbox, y referencias de pago válidas.
  \item La documentación indica claramente la fuente oficial de cada dato (enlace a la documentación del proveedor).
  \item La documentación se mantiene actualizada con los datos de prueba vigentes de cada proveedor al momento de la entrega del proyecto.
\end{itemize}